%====================================================================
%====================================================================
\section*{Graphical models}
%====================================================================

\renewcommand{\pause}{}

%====================================================================
\subsection*{Directed graphical models}
%==================================================================
\frame{\frametitle{A simple (interesting) example}

  $p$ faithful to $D =$
  $$
  \input{\figeco/LeftRightChainABC}
  $$
  means that
  $$
  p(a, b, c) = p(a) \; p(b \mid a) \; p(c \mid a) 
  $$ \pause

  \bigskip \bigskip
  $p$ is also faithful to $D' =$
  $$
  \input{\figeco/RightLeftChainABC} 
  $$
  \pause and to $D'' =$
  $$
  \input{\figeco/AwayFromCenterABC}
  $$
  \pause but not to $D''' =$
  $$
  \input{\figeco/HeadToHeadABC}
  $$  
}
  
%====================================================================
\frame{\frametitle{Interpretability?}

  \paragraph{Theorem \refer{Pea09}.} 
  Two DAGs are Markov equivalent (i.e. induce the same conditional dependences and independences) if they have the \emphase{same skeleton} (i.e. the same undirected edges) and the \emphase{same V-structures}.

  \bigskip \bigskip \pause
  \paragraph{Example.}
  \begin{center}  
  \begin{tabular}{c|cc|cc}
  $D$ &
  \multicolumn{2}{c|}{Equivalent to $D$} &
  \multicolumn{2}{c}{Not equivalent to $D$} \\
  (\textcolor{red}{V-stuctures}) & & & & \\ 
%   & & & & \\ \hline & & & & \\ 
  \hline
  $\footnotesize{\input{\figeco/DAG6nodes-vtruct}}$ \pause &
  $\footnotesize{\input{\figeco/DAG6nodes-vtruct-equiv1}}$ &
  $\footnotesize{\input{\figeco/DAG6nodes-vtruct-equiv2}}$ \pause &
  $\footnotesize{\input{\figeco/DAG6nodes-vtruct-notequiv1}}$ &
  $\footnotesize{\input{\figeco/DAG6nodes-vtruct-notequiv2}}$ 
  \end{tabular}
  \end{center}
  (Not speaking of possibly missing nodes)
}

%==================================================================
\subsection*{Gaussian graphical models}
%==================================================================
\frame{\frametitle{Regression point of view}

  \paragraph{Similarly to temporal networks,} finding the neighbors of a node is equivalent to find its 'parents' in a regression:
  $$
  Y_{ik} = \sum_{j \neq k} \beta_{jk} Y_{ij} + E_{ik}.
  $$
  \ra Requires a testing procedure \refer{VeV09} or a penalization \refer{MeB06} to determine non-zero coefficients

  \bigskip \bigskip \pause
  \paragraph{Remarks.}
  \begin{itemize}
   \item Regression needs reconciliation when $\widehat{\beta}_{jk} = 0$ and $\widehat{\beta}_{jk} \neq 0$ 
   \item \refer{Ver12} provides bounds for the recovery of the list of neighbors: let $k =$ degree of a given node 
   \begin{align*}
    k \log(p/k) & \leq n & & \text{possible recovery} \\
    k \log(p/k) & > n \log n & & \text{impossible recovery}
   \end{align*}
  \end{itemize}
}

%==================================================================
\frame{\frametitle{GGM's nice property: More formally}

  \paragraph{If $Y \sim \Ncal(0, \Sigma)$,} then
  \begin{align*}
  p(Y) & \propto \exp\left(-\frac12 \|Y\|^2_{\Sigma^{-1}} \right)\\
  & = \exp\left(-\frac12 \sum_{j, k} \omega_{jk} Y_j Y_k\right)  
  & & \text{where } \Omega = [\omega_{jk}] = \Sigma^{-1} \\
  & = \prod_{j, k} \underset{\phi_{jk}(Y_j, Y_k)}{\underbrace{\exp\left(-\frac12 \omega_{jk} Y_j Y_k\right)}}
  \end{align*}

  \bigskip \pause
  \begin{itemize}
   \item The non-zeros of $\Omega$ correspond to the edges of $G$
   \item Furthermore:
   $$
   - \omega_{jk} \; \propto \; 
   \rho\left(Y_j, Y_k \mid Y_{\{j, k\}}\right)
   = \text{ \emphase{'partial'} correlation}
   $$
  \end{itemize}
  
}


%====================================================================
%====================================================================
\section*{Network inference from count data}
%====================================================================

%====================================================================
%====================================================================
\section*{Extensions}
%====================================================================

%====================================================================
\subsection*{Missing actors}
%====================================================================
\frame{\frametitle{Missing actors}
}
